\documentclass[10pt,a4paper]{amsart}
\usepackage{amsmath,amssymb,amsthm,graphicx}
% Retained minimal preamble for the Stacks Project chapter coherent.tex (Stacks chapter, pre-delegation retained fragment for stacks-coherent).
% Verbatim excerpt of the chapter preamble (sources/stacks-spaces-descent/preamble.tex): the reference/slogan/history
% environment block (native lines 7-12), Xy-pic setup (native lines 14-19) and the native theorem declarations with
% the mathematical macros (native lines 154-217). Copied mechanically; nothing added, removed or redefined. The
% wrapper class and the amsmath/amssymb/amsthm/graphicx base packages are supplied by the assembler, so the
% class/xr-hyper/hyperref/externaldocument part of the native preamble is omitted.
% For dealing with references we use the comment environment
\usepackage{verbatim}
\newenvironment{reference}{\comment}{\endcomment}
%\newenvironment{reference}{}{}
\newenvironment{slogan}{\comment}{\endcomment}
\newenvironment{history}{\comment}{\endcomment}
% For commutative diagrams we use Xy-pic
\usepackage[all]{xy}

% We use 2cell for 2-commutative diagrams.
\xyoption{2cell}
\UseAllTwocells
% Theorem environments.
%
\theoremstyle{plain}
\newtheorem{theorem}[subsection]{Theorem}
\newtheorem{proposition}[subsection]{Proposition}
\newtheorem{lemma}[subsection]{Lemma}

\theoremstyle{definition}
\newtheorem{definition}[subsection]{Definition}
\newtheorem{example}[subsection]{Example}
\newtheorem{exercise}[subsection]{Exercise}
\newtheorem{situation}[subsection]{Situation}

\theoremstyle{remark}
\newtheorem{remark}[subsection]{Remark}
\newtheorem{remarks}[subsection]{Remarks}

\numberwithin{equation}{subsection}

% Macros
%
\def\lim{\mathop{\mathrm{lim}}\nolimits}
\def\colim{\mathop{\mathrm{colim}}\nolimits}
\def\Spec{\mathop{\mathrm{Spec}}}
\def\Hom{\mathop{\mathrm{Hom}}\nolimits}
\def\Ext{\mathop{\mathrm{Ext}}\nolimits}
\def\SheafHom{\mathop{\mathcal{H}\!\mathit{om}}\nolimits}
\def\SheafExt{\mathop{\mathcal{E}\!\mathit{xt}}\nolimits}
\def\Sch{\mathit{Sch}}
\def\Mor{\mathop{\mathrm{Mor}}\nolimits}
\def\Ob{\mathop{\mathrm{Ob}}\nolimits}
\def\Sh{\mathop{\mathit{Sh}}\nolimits}
\def\NL{\mathop{N\!L}\nolimits}
\def\CH{\mathop{\mathrm{CH}}\nolimits}
\def\proetale{{pro\text{-}\acute{e}tale}}
\def\etale{{\acute{e}tale}}
\def\QCoh{\mathit{QCoh}}
\def\Ker{\mathop{\mathrm{Ker}}}
\def\Im{\mathop{\mathrm{Im}}}
\def\Coker{\mathop{\mathrm{Coker}}}
\def\Coim{\mathop{\mathrm{Coim}}}

% Boxtimes
%
\DeclareMathSymbol{\boxtimes}{\mathbin}{AMSa}{"02}

%
% Macros for moduli stacks/spaces
%
\def\QCohstack{\mathcal{QC}\!\mathit{oh}}
\def\Cohstack{\mathcal{C}\!\mathit{oh}}
\def\Spacesstack{\mathcal{S}\!\mathit{paces}}
\def\Quotfunctor{\mathrm{Quot}}
\def\Hilbfunctor{\mathrm{Hilb}}
\def\Curvesstack{\mathcal{C}\!\mathit{urves}}
\def\Polarizedstack{\mathcal{P}\!\mathit{olarized}}
\def\Complexesstack{\mathcal{C}\!\mathit{omplexes}}
% \Pic is the operator that assigns to X its picard group, usage \Pic(X)
% \Picardstack_{X/B} denotes the Picard stack of X over B
% \Picardfunctor_{X/B} denotes the Picard functor of X over B
\def\Pic{\mathop{\mathrm{Pic}}\nolimits}
\def\Picardstack{\mathcal{P}\!\mathit{ic}}
\def\Picardfunctor{\mathrm{Pic}}
\def\Deformationcategory{\mathcal{D}\!\mathit{ef}}
\usepackage{cleveref}
\crefname{originallemma}{original lemma}{original lemmas}
\Crefname{originallemma}{Original lemma}{Original lemmas}
\makeatletter
\newlabel{lemma-coherent-filter}{{01YF}{source}{}{sourcepointer-1}{}}
\newlabel{lemma-coherent-filter@cref}{{[originallemma][0][]01YF}{[1][0][]source}{}{}{}}
\makeatother
\crefname{originallemma}{original lemma}{original lemmas}
\Crefname{originallemma}{Original lemma}{Original lemmas}
\makeatletter
\newlabel{lemma-power-ideal-kills-sheaf}{{01Y9}{source}{}{sourcepointer-2}{}}
\newlabel{lemma-power-ideal-kills-sheaf@cref}{{[originallemma][0][]01Y9}{[1][0][]source}{}{}{}}
\makeatother
\crefname{originallemma}{original lemma}{original lemmas}
\Crefname{originallemma}{Original lemma}{Original lemmas}
\makeatletter
\newlabel{lemma-prepare-filter-irreducible}{{01YE}{source}{}{sourcepointer-3}{}}
\newlabel{lemma-prepare-filter-irreducible@cref}{{[originallemma][0][]01YE}{[1][0][]source}{}{}{}}
\makeatother
\title{Devissage with higher rank on an irreducible closed subset: for X Noetherian, Z\_0 irreducible closed with generic point xi, a property P of coherent sheaves stable under extensions, descending from P(F\textasciicircum{}\{oplus r\}) to P(F), holding for pushforwards of ideal sheaves along integral closed subschemes properly inside Z\_0, and for some coherent G supported on Z\_0 with G\_xi annihilated by m\_xi such that every ideal J with J\_xi = O\_\{X,xi\} admits a subsheaf G' of JG with G'\_xi = G\_xi satisfying P: then P holds for every coherent sheaf whose support is contained in Z\_0.}
\author{The Stacks Project Authors}
\date{stacks-project:coherent.tex:local-snapshot-20261012; GNU FDL 1.2 or later}
\begin{document}
\maketitle
\noindent\textit{Editorial note.} Statement span is lines 3523-3555 of sources/stacks-coherent/pinned-original-source.tex; proof span is lines 3557-3639, the complete author proof (81 body lines).\par\medskip
\noindent\textit{Editorial note.} Navigation stubs: 3 cited-but-not-included labels are supplied as non-load-bearing pointers with their pinned tags where the snapshot tag index has them.\par\medskip
\noindent\textit{Editorial note.} extra\_packages is empty on purpose: the retained preamble already loads xy with options and a second declaration clashes in the assembler wrapper.\par\medskip
\noindent\textit{Editorial note.} Text is retained verbatim, including the author's own typos and placeholders.\par\medskip
\noindent\textit{Editorial note.} Full derivation in span; no omissions.\par\medskip
\noindent\textit{Editorial note.} Span computed from the \textbackslash{}label marker through the enclosing \textbackslash{}end\{lemma\}; precard\_gate BLOCK was the public-interval false overlap described in dedup (semantic occupancy verified manually).\par\medskip
\noindent\textit{Allowed standard prerequisites (editorial).}\begin{itemize}
\item Explicit prerequisite boundary. See why-field; every pointer named here occurs verbatim in the retained statement/proof text.
\end{itemize}\medskip
\ifdefined\hypertarget\hypertarget{sourcepointer-1}{}\fi
\noindent\textit{Source reference.} Original lemma 01YF: sources/stacks-coherent/pinned-original-source.tex, source lines 3268--3284. Cited by the retained statement or proof but not part of the collected spans; supplied only as a navigation pointer. The argument does not depend on its content.\par\medskip
\ifdefined\hypertarget\hypertarget{sourcepointer-2}{}\fi
\noindent\textit{Source reference.} Original lemma 01Y9: sources/stacks-coherent/pinned-original-source.tex, source lines 2786--2793. Cited by the retained statement or proof but not part of the collected spans; supplied only as a navigation pointer. The argument does not depend on its content.\par\medskip
\ifdefined\hypertarget\hypertarget{sourcepointer-3}{}\fi
\noindent\textit{Source reference.} Original lemma 01YE: sources/stacks-coherent/pinned-original-source.tex, source lines 3199--3212. Cited by the retained statement or proof but not part of the collected spans; supplied only as a navigation pointer. The argument does not depend on its content.\par\medskip
\begin{lemma}
\label{lemma-property-irreducible-higher-rank-cohomological}
Let $X$ be a Noetherian scheme. Let $Z_0 \subset X$ be an irreducible
closed subset with generic point $\xi$. Let $\mathcal{P}$ be a property
of coherent sheaves on $X$ such that
\begin{enumerate}
\item For any short exact sequence of coherent sheaves
$$
0 \to \mathcal{F}_1 \to \mathcal{F} \to \mathcal{F}_2 \to 0
$$
if $\mathcal{F}_i$, $i = 1, 2$ have property $\mathcal{P}$
then so does $\mathcal{F}$.
\item If $\mathcal{P}$ holds for $\mathcal{F}^{\oplus r}$ for
some $r \geq 1$, then it holds for $\mathcal{F}$.
\item For every integral closed subscheme $Z \subset Z_0 \subset X$,
$Z \not = Z_0$ and every quasi-coherent sheaf of ideals
$\mathcal{I} \subset \mathcal{O}_Z$ we have
$\mathcal{P}$ for $(Z \to X)_*\mathcal{I}$.
\item There exists some coherent sheaf $\mathcal{G}$ such that
\begin{enumerate}
\item $\text{Supp}(\mathcal{G}) = Z_0$,
\item $\mathcal{G}_\xi$ is annihilated by $\mathfrak m_\xi$, and
\item for every quasi-coherent sheaf of ideals
$\mathcal{J} \subset \mathcal{O}_X$ such that
$\mathcal{J}_\xi = \mathcal{O}_{X, \xi}$ there exists a quasi-coherent
subsheaf $\mathcal{G}' \subset \mathcal{J}\mathcal{G}$ with
$\mathcal{G}'_\xi = \mathcal{G}_\xi$ and such that
$\mathcal{P}$ holds for $\mathcal{G}'$.
\end{enumerate}
\end{enumerate}
Then property $\mathcal{P}$ holds for every coherent sheaf
$\mathcal{F}$ on $X$ whose support is contained in $Z_0$.
\end{lemma}

\begin{proof}
Note that if $\mathcal{F}$ is a coherent sheaf with a filtration
$$
0 = \mathcal{F}_0 \subset \mathcal{F}_1 \subset
\ldots \subset \mathcal{F}_m = \mathcal{F}
$$
by coherent subsheaves such that each of $\mathcal{F}_i/\mathcal{F}_{i - 1}$
has property $\mathcal{P}$, then so does $\mathcal{F}$.
This follows from assumption (1).

\medskip\noindent
As a first application we conclude that any coherent sheaf whose support
is strictly contained in $Z_0$ has property $\mathcal{P}$. Namely, such a
sheaf has a filtration (see Lemma \ref{lemma-coherent-filter})
whose subquotients have property $\mathcal{P}$ according to (3).

\medskip\noindent
Let us denote $i : Z_0 \to X$ the closed immersion.
Consider a coherent sheaf $\mathcal{G}$ as in (4).
By Lemma \ref{lemma-prepare-filter-irreducible}
there exists a sheaf of ideals $\mathcal{I}$ on $Z_0$ and
a short exact sequence
$$
0 \to
i_*\mathcal{I}^{\oplus r} \to
\mathcal{G} \to
\mathcal{Q} \to 0
$$
where the support of $\mathcal{Q}$ is strictly contained in $Z_0$.
In particular $r > 0$ and $\mathcal{I}$ is nonzero
because the support of $\mathcal{G}$ is equal to $Z_0$.
Let $\mathcal{I}' \subset \mathcal{I}$ be any nonzero quasi-coherent
sheaf of ideals on $Z_0$ contained in $\mathcal{I}$.
Then we also get a short exact sequence
$$
0 \to
i_*(\mathcal{I}')^{\oplus r} \to
\mathcal{G} \to
\mathcal{Q}' \to 0
$$
where $\mathcal{Q}'$ has support properly contained in $Z_0$.
Let $\mathcal{J} \subset \mathcal{O}_X$ be a quasi-coherent sheaf
of ideals cutting out the support of $\mathcal{Q}'$ (for example
the ideal corresponding to the reduced induced closed subscheme
structure on the support of $\mathcal{Q}'$). Then
$\mathcal{J}_\xi = \mathcal{O}_{X, \xi}$. By
Lemma \ref{lemma-power-ideal-kills-sheaf}
we see that $\mathcal{J}^n\mathcal{Q}' = 0$ for some $n$.
Hence $\mathcal{J}^n\mathcal{G} \subset i_*(\mathcal{I}')^{\oplus r}$.
By assumption (4)(c) of the lemma we see there exists
a quasi-coherent subsheaf $\mathcal{G}' \subset \mathcal{J}^n\mathcal{G}$
with $\mathcal{G}'_\xi = \mathcal{G}_\xi$
for which property $\mathcal{P}$ holds.
Hence we get a short exact sequence
$$
0 \to \mathcal{G}' \to
i_*(\mathcal{I}')^{\oplus r} \to
\mathcal{Q}'' \to 0
$$
where $\mathcal{Q}''$ has support properly contained in $Z_0$.
Thus by our initial remarks and property (1) of the lemma
we conclude that $i_*(\mathcal{I}')^{\oplus r}$ satisfies
$\mathcal{P}$. Hence we see that $i_*\mathcal{I}'$ satisfies
$\mathcal{P}$ by (2). Finally, for an arbitrary quasi-coherent
sheaf of ideals $\mathcal{I}'' \subset \mathcal{O}_{Z_0}$ we can set
$\mathcal{I}' = \mathcal{I}'' \cap \mathcal{I}$ and we get
a short exact sequence
$$
0 \to
i_*(\mathcal{I}') \to
i_*(\mathcal{I}'') \to
\mathcal{Q}''' \to 0
$$
where $\mathcal{Q}'''$ has support properly contained in $Z_0$.
Hence we conclude that property $\mathcal{P}$ holds for
$i_*\mathcal{I}''$ as well.

\medskip\noindent
The final step of the proof is to note that any coherent sheaf
$\mathcal{F}$ on $X$ whose support is contained in $Z_0$ has a filtration
(see Lemma \ref{lemma-coherent-filter} again) whose subquotients
all have property $\mathcal{P}$ by what we just said.
\end{proof}

\end{document}
